\documentclass{article}
\usepackage{amsmath, amssymb, amsthm}
\usepackage{mathpazo}

\newtheorem{theorem}{Theorem}
\newtheorem{lemma}{Lemma}

\begin{document}

\section*{Formal Specification: JQ $\to$ SQL Correctness}

This document formalizes a translation from jq-style queries over JSON data into SQL-style queries over relational data, and proves that the translation preserves meaning.

% =====================================================

\subsection*{1. Data Structures and Domains}

This section defines the two representations of the same logical data: a user database stored as JSON records and as SQL tuples.

Let $\mathcal{J}$ be the domain of JSON users, where each $j \in \mathcal{J}$ is a record:
\[
j = \{ \text{name} : \text{String}, \ \text{age} : \mathbb{N} \}.
\]

Let $\mathcal{S}$ be the domain of SQL users, where each $s \in \mathcal{S}$ is a tuple:
\[
s = (\text{name}, \text{age}) \in \text{String} \times \mathbb{N}.
\]

\begin{itemize}
    \item \textbf{JSON Database ($\mathcal{JD}$):} A finite list of JSON user records.
    \item \textbf{SQL Database ($\mathcal{SD}$):} A finite list of SQL tuples.
\end{itemize}

We define conversion functions that preserve structure:
\[
\text{toS} : \mathcal{J} \to \mathcal{S}, \quad
\text{toJ} : \mathcal{S} \to \mathcal{J}
\] 
such that:
\[ \text{toS}(\{\text{name}=n, \text{age}=a\}) = (n, a) \]

These functions convert between equivalent representations of the same data.

% =====================================================

\subsection*{2. Database Equivalence}

This section defines when a JSON database and an SQL database represent the same information.

Two databases $JD \in \mathcal{JD}$ and $SD \in \mathcal{SD}$ are equivalent, written $JD \equiv SD$, if they contain corresponding data under conversion:

\[
JD \equiv SD \iff
\left( \forall j \in \mathcal{J},\ j \in JD \leftrightarrow \text{toS}(j) \in SD \right)
\land
\left( \forall s \in \mathcal{S},\ s \in SD \leftrightarrow \text{toJ}(s) \in JD \right).
\]

Intuitively, every JSON record exists in the database if and only if its SQL encoding exists in the SQL database, and vice versa.

% =====================================================

\subsection*{3. Condition Semantics}

This section defines how conditions (e.g. \texttt{age > 20}) are evaluated in both representations.

A condition $C$ is interpreted over both formats:

\begin{itemize}
    \item $\text{eval}_J(C, j)$: evaluates condition $C$ on JSON record $j$
    \item $\text{eval}_S(C, s)$: evaluates condition $C$ on SQL tuple $s$
\end{itemize}

\paragraph{Bridge Lemma}
This lemma ensures that condition evaluation is consistent across representations.

\begin{equation}
\forall C,\ \forall j \in \mathcal{J},\quad
\text{eval}_J(C, j) = \text{eval}_S(C, \text{toS}(j)).
\end{equation}

This is the key semantic consistency property between JSON and SQL worlds.

% =====================================================

\subsection*{4. Translation Function}

This section defines the mapping from jq queries to SQL queries.

\[
T : Q_{jq} \to Q_{sql}
\]

The translation preserves the structure of the query while changing its target representation:

\begin{itemize}
    \item \texttt{find} $\rightarrow$ \texttt{SELECT}
    \item \texttt{drop} $\rightarrow$ \texttt{DELETE}
\end{itemize}

The condition $C$ remains unchanged during translation.

% =====================================================

\subsection*{5. Query Semantics}

This section defines how queries operate as filters over databases.

A query $q$ selects or deletes elements based on a condition $C$.

\begin{itemize}
\item \textbf{jq semantics:} returns all JSON records satisfying $C$
\[
\text{eval}_{jq}(JD, q) =
\{ j \in JD \mid \text{eval}_J(C, j) = \text{true} \}
\]

\item \textbf{SQL semantics:} returns all SQL tuples satisfying $C$
\[
\text{eval}_{sql}(SD, T(q)) =
\{ s \in SD \mid \text{eval}_S(C, s) = \text{true} \}
\]
\end{itemize}

Both systems behave as filters over their respective databases.

% =====================================================

\subsection*{6. Main Correctness Theorem}

This theorem states that translation preserves meaning across both database systems.

\begin{theorem}
Let $JD \in \mathcal{JD}$ and $SD \in \mathcal{SD}$ such that $JD \equiv SD$.  
For any jq query $q$, execution after translation preserves equivalence:
\[
\text{eval}_{jq}(JD, q)
\equiv
\text{toS}\big(\text{eval}_{sql}(SD, T(q))\big).
\]
\end{theorem}

\begin{proof}
We prove element-wise equivalence.

For any $j$:
\[
j \in \text{eval}_{jq}(JD, q)
\iff
j \in JD \land \text{eval}_J(C,j)
\]

Using database equivalence:
\[
j \in JD \iff \text{toS}(j) \in SD
\]

Using the Bridge Lemma:
\[
\text{eval}_J(C,j) \iff \text{eval}_S(C,\text{toS}(j))
\]

Combining both:
\[
j \in \text{eval}_{jq}(JD,q)
\iff
\text{toS}(j) \in \text{eval}_{sql}(SD,T(q))
\]

Hence the result follows.
\end{proof}

% =====================================================

\subsection*{7. Interpretation}

This result guarantees semantic preservation: executing a jq query on JSON data produces results equivalent (after conversion) to executing the translated SQL query on relational data.
In other words, the translation function $T$ is semantics-preserving.

\end{document}